\documentclass[11pt]{article}

\usepackage[T1]{fontenc}
\usepackage{lmodern}
% Leave room in the right margin for marginalia and todo notes.
\usepackage[margin=1in,marginparwidth=1.5in,marginparsep=0.2in]{geometry}
\usepackage{amsmath}
\usepackage{amssymb}
\usepackage{marginnote}
\usepackage{enumitem}
\usepackage{todonotes}
\usepackage{cleveref}

\setlength{\marginparwidth}{1.5in}

\title{Marginalia and Rich Footnotes as a Conversion Stress Case}
\author{Test Corpus}
\date{}

\begin{document}
\maketitle

\begin{abstract}
This fixture concentrates the apparatus that academic writing pushes into the
margins and the foot of the page: margin notes set with both \texttt{marginnote}
and \texttt{marginpar}, a \texttt{todonotes} annotation, and footnotes that carry
inline mathematics, a cross-reference, and an enumerated list. Each construct
detaches content from the main text flow, so a faithful conversion must relocate
that content correctly rather than dropping it or splicing it inline.
\end{abstract}

\section{Margin Notes}
\label{sec:margins}

Margin notes annotate the body text without interrupting it.\marginnote{Set with
\texttt{marginnote}; unlike \texttt{marginpar} it does not float and may be used
inside floats and footnotes.} The note beside this sentence is produced with
\verb|\marginnote|, which places fixed text in the adjacent margin and, because it
does not float, behaves predictably even in contexts where a floating note would
fail. A second mechanism, \verb|\marginpar|, lets \LaTeX{} choose the margin and
shift the note vertically to avoid collisions.\marginpar{\footnotesize A
\texttt{marginpar} note, allowed to float to avoid overlap.} Both mechanisms
share the property that their content is anchored to a point in the text but
typeset outside the main column, which is precisely the relocation a converter
must reproduce.

\todo[inline]{Confirm that the inline todo note survives conversion as a visible,
set-apart annotation rather than vanishing or merging into the body paragraph.}

The annotation above is a \verb|\todo| note from the \texttt{todonotes} package.
In the compiled article it is rendered as a colored, set-apart box, and it stands
in for the editorial markers that pervade working drafts. \Cref{sec:footnotes}
turns from the margins to the foot of the page.

\section{Rich Footnotes}
\label{sec:footnotes}

Footnotes are the second channel by which scholarly prose offloads material from
the main flow, and they frequently carry more than plain text. The footnote
attached to this sentence contains inline mathematics and a cross-reference back
into the body.\footnote{The estimator is unbiased, so its mean equals the target,
$\mathbb{E}[\hat{\theta}] = \theta$, with sampling variance
$\operatorname{Var}(\hat{\theta}) = \sigma^{2}/n$; this expands on the model named
in \cref{sec:margins}.} Inline mathematics inside a footnote is a common source of
conversion error, because the equation must remain mathematics even after the note
is moved to its own region at the foot of the page.

A footnote may also enumerate several items rather than running them together in a
sentence.\footnote{Three conditions support the variance expression: %
\begin{enumerate}[label=(\roman*),nosep,leftmargin=*]
  \item the disturbances are uncorrelated across observations;
  \item they share a common variance $\sigma^{2}$; and
  \item the regressors are fixed in repeated sampling.
\end{enumerate}}
The list nested in that footnote tests whether enumerated structure is preserved
when it appears in a footnote rather than in the body, a placement that some
converters silently flatten.

\section{Interaction}
\label{sec:interaction}

The preceding sections isolate each construct; this one places them in proximity
to mirror a dense passage of real scholarship.\marginnote{\footnotesize Margin
notes and footnotes routinely coexist in the same paragraph.} A claim in the body
can carry a margin gloss while simultaneously anchoring a footnote that qualifies
it,\footnote{Footnotes and margin notes use independent numbering and placement,
so a converter must not conflate the two even when they originate from adjacent
points in the source.} and an editorial \verb|\todo| may sit nearby to flag
unfinished work.

\todo[inline]{Verify that a second todo note, placed after the footnote-heavy
paragraph, is still rendered distinctly and is not absorbed by the surrounding
text.}

Taken together, \cref{sec:margins,sec:footnotes,sec:interaction} establish that
the document exercises marginalia, editorial annotations, and footnotes that
themselves contain mathematics, cross-references, and lists. A faithful
conversion must preserve each of these out-of-flow elements and keep the
cross-references resolvable across the sections that house them.

\end{document}
